% This template was provided by Dr. Charles Welch
% He uses it in his course
%

\documentclass[11pt]{article}

% Change "review" to "final" to generate the final (sometimes called camera-ready) version.
% Change to "preprint" to generate a non-anonymous version with page numbers.
\usepackage[]{acl}
\usepackage{times}
\usepackage{latexsym}

% For proper rendering and hyphenation of words containing Latin characters (including in bib files)
\usepackage[T1]{fontenc}
% For Vietnamese characters
% \usepackage[T5]{fontenc}
% See https://www.latex-project.org/help/documentation/encguide.pdf for other character sets

% This assumes your files are encoded as UTF8
\usepackage[utf8]{inputenc}
\usepackage{microtype}
\usepackage{inconsolata}
\usepackage{graphicx}
\usepackage{booktabs}

% If the title and author information does not fit in the area allocated, uncomment the following
%
%\setlength\titlebox{<dim>}
%
% and set <dim> to something 5cm or larger.

\title{ML Report:\\Bridging Gaps: AI for Diagram Accessibility}


\author{Nawaal Fatima, Casey Francine Bulaclac, Moly Mikhail, Fiza Sehar, Dhruv Sardana \\
  \texttt{\{fatimn8,bulacc1,mikham16,seharf,sardanad\}@mcmaster.ca} }

\begin{document}
\maketitle

\section{Introduction}
This report outlines the development of the artificial intelligence component of the Reading4All capstone project, which focuses on generating descriptive alternative text for physics diagrams to 
improve accessibility for visually impaired users.

To achieve this goal, we explored and evaluated multiple approaches, going through several iterations and experiments to refine our model's performance and output quality. Early stages involved 
testing different model architectures, preprocessing techniques, and training strategies to better understand how to generate accurate and readable descriptions for technical diagrams.

Through this iterative process, we identified key limitations in initial outputs, such as lack of detail, inconsistency, and formatting issues. These insights guided further experimentation 
and improvements, including dataset refinement, training adjustments, and post-processing enhancements.

The final model selected for this project is based on BLIP (Bootstrapping Language-Image Pretraining), which provided the best balance between descriptive accuracy and generalization. 
This report details the progression throughout the different revisions from initial experimentation to the finalized model, highlighting the design decisions, challenges, and improvements made along the way.

\section{Related Work}

This section discusses existing work most relevant to our approach, focusing on vision-language models used for image captioning and structured diagram understanding. While no prior work addresses the specific task of generating alt text for physics diagrams, the models and techniques described below represent the closest related work and directly informed our design decisions.

\subsection{Vision-Language Models}
\subsubsection{BLIP}
Bootstrapping Language-Image Pretraining (BLIP) is a vision-language model designed for tasks that require joint understanding of images and text, such as image captioning and visual question answering. 
It is particularly relevant to this project because it generates natural language descriptions from visual inputs, making it well-suited for automatic alt-text generation.

BLIP uses a unified architecture that combines a vision encoder with a language model, allowing it to handle both image-text understanding and text generation. It is trained using multiple objectives, 
including image-text contrastive learning, matching, and language modeling, which help it learn strong and generalizable representations.
In this project, BLIP was selected after multiple iterations and experiments due to its ability to generate more coherent, detailed, and structured alt text for complex physics diagrams. It also 
supports fine-tuning and post-processing, which further improves readability and consistency.

BLIP provides a strong and flexible foundation for generating accessible image descriptions, aligning well with the goals of this project.

\subsubsection{Pix2Struct}
Pix2Struct is a visually-situated language model designed for parsing and understanding 
structured visual inputs such as charts, diagrams, and document screenshots. Unlike 
general-purpose vision-language models, Pix2Struct is pretrained on screenshot-to-text 
tasks, making it particularly relevant for structured scientific figures. It processes 
images by rendering them as a sequence of patches and learning to generate textual 
descriptions directly from visual layout. This design makes it a strong candidate for 
diagram-heavy domains such as physics, where spatial relationships and annotations are 
critical to understanding. In this project, Pix2Struct was evaluated as an alternative 
to BLIP, but ultimately underperformed on our custom physics dataset due to slower 
inference, a tendency to produce overly generic or truncated captions, and difficulty 
generalizing to domain-specific diagrams outside its pretraining distribution.

\section{Dataset}

This section describes the dataset used to train and evaluate the model, including its composition, annotation structure, and preprocessing steps applied across revisions.

\subsection{Dataset Overview}

The dataset focused exclusively on physics-related diagrams to align with the goals of the project. The domains covered included topics such as magnetic fields, mechanics, and various types of graphs. By limiting the dataset to 
these specific areas, the model was better able to learn relevant patterns and produce more precise and context-aware alt text for physics content.

In Revision 0, the initial dataset consisted of approximately 500 images. These images were primarily physics diagrams used to establish a baseline for training and evaluating the model. In Revision 1, the dataset was expanded 
by adding approximately 300 additional images, bringing the total to around 800 images. This expansion helped improve the model's ability to generalize and generate more accurate descriptions.

\subsection{Annotation Structure}
In Revision 0, the dataset annotations were stored in a CSV file with four main fields: Diagram-ID, Image-Path, Category-of-Image, and Modified-Alt-Text. The Diagram-ID serves as 
a unique identifier for each image, allowing for easy tracking and referencing. The Image-Path specifies the file location of the corresponding diagram. The Modified-Alt-Text contains 
the descriptive text associated with each image, 
which serves as the ground truth for training the model. These descriptions are detailed and structured to capture key components, relationships, and annotations present in the diagrams.

In Revision 1, an additional field called Category-of-Image was introduced to further refine the categorization of the dataset. This field groups diagrams into more specific types, such as graphs, 
waves, mechanics, and others. Adding this column allowed for better control and specificity in prompting the model, ultimately improving the relevance and consistency of the generated alt text.

\subsection{Preprocessing}

Minimal preprocessing was applied to preserve the original structure and details of the physics diagrams. Images were not manually resized or normalized, as maintaining the original resolution helped retain fine-grained elements such as labels, arrows, and annotations that are important for accurate alt-text generation. However, the BLIP processor automatically handled resizing and normalization internally to match the model's expected input format during training.

The primary preprocessing effort focused on cleaning and refining the alt-text annotations. Existing alt text descriptions were reviewed and reformatted to improve consistency, readability, and structure. This included correcting formatting issues, standardizing phrasing, and ensuring that descriptions clearly captured relationships between diagram components.

User feedback played an important role in this process. Approximately 150 images had previously written alt text that was evaluated by users. Based on this feedback, the alt text was modified to improve clarity, accessibility, and descriptive quality. In addition to these refined annotations, approximately 700 additional images were added to the dataset with alt text generated from scratch. These new descriptions were written following the same structured format to maintain consistency across the dataset.

Additionally, the dataset was filtered to remove duplicate or low-quality images that could negatively impact training. This helped maintain dataset consistency and reduced noise during model learning.

Finally, the dataset was split into training, validation, and testing sets to support structured experimentation and evaluation. This split allowed for performance monitoring while preventing overfitting and ensuring that the model was evaluated on unseen data.

\section{Features}

This section describes the features and input representations used across model iterations, including how data was fed into each model and what forms of feature learning were applied. It covers the baseline setup in Revision 0 and the series of experiments conducted in Revision 1 to improve output quality.

\subsection{Baseline (Revision 0)}
In Revision 0, we experimented with either using BLIP or Pix2Struct for our AI model. 

\subsubsection{BLIP}
In Revision 0, BLIP was used as the baseline model for generating alt text from physics diagrams. The model takes an 
image as input and processes it through a vision encoder to extract visual features such as shapes, text, and spatial 
relationships within the diagram. These features are then passed to a language model, which generates a natural 
language description of the image. This process relies on learned embeddings, where both visual and textual information 
are represented in a shared feature space, allowing the model to effectively connect image content with descriptive text.

No explicit feature engineering or manual feature selection was performed, as BLIP inherently learns relevant features 
during training through its pretraining objectives. The model was used in a zero-shot or minimally fine-tuned setting 
in Revision 0, meaning it relied primarily on its pretrained knowledge to generate descriptions. While this allowed for 
quick implementation, the outputs often lacked consistency, structure, and domain-specific detail for complex physics diagrams.

BLIP in Revision 0 served as a strong starting point by providing automatic caption generation, but it highlighted 
the need for further improvements in dataset quality, domain specificity, and output formatting in later revisions.

\subsubsection{Pix2Struct}
In Revision 0, Pix2Struct was also evaluated as a potential model for alt text generation. 
Like BLIP, it takes an image as input and generates a natural language description, but 
it differs in that it is specifically pretrained on structured visual inputs such as 
charts, diagrams, and web screenshots. Pix2Struct processes images by dividing them into 
a sequence of visual patches, which are then used to generate text through a 
transformer-based decoder. No manual feature engineering was applied; the model relied 
entirely on its pretrained representations. However, in practice, Pix2Struct struggled 
to generalize to our custom physics diagrams, producing descriptions that were overly 
brief, sometimes hallucinated content, and were slow to run given our hardware constraints. 
These limitations led to BLIP being selected as the stronger baseline for subsequent revisions.

\subsection{Revision 1}
In Revision 1, we conducted different experiments highlighted below to see if we can improve our 
AI model from Revision 0.

\subsubsection{BLIP Re-training for Expanded Dataset} 
In Revision 1, we aimed to improve our model by expanding our training dataset. 
As new images were annotated, we iteratively retrained the model using larger datasets
(100, 150 and 200 new images). This allowed us to compare the model's output on 
test images and observe gradual improvements in performance. 
Through this process, we were able to experiment with different dataset sizes
and evaluate the trade-off between additional data and performance increases. 

\subsubsection{Category Routing System}

To improve the relevance of generated alt text, a category routing system was introduced in Revision 1. The dataset contained multiple diagram categories representing different physics concepts, and a single captioning model often struggled to generate accurate descriptions across all diagram types.

To address this, a CLIP-based classifier was trained to categorize each input diagram before caption generation. The goal of this approach was to first identify the diagram type using the CLIP model and then use this category information to guide the prompt provided to the BLIP model. By conditioning the prompt on the predicted category, the system was able to generate more context-aware and structured alt text tailored to each diagram type.

Additionally, a subject field was incorporated into the dataset to provide further contextual guidance. This subject information, combined with the predicted category, helped refine the prompting strategy and improved consistency when generating descriptions for visually similar diagrams belonging to different conceptual domains.

This routing strategy allowed the system to leverage CLIP for category prediction and BLIP for caption generation, improving descriptive accuracy and reducing overly generic outputs.

\subsubsection{LLM Polisher Post-Processing}

In Revision 1, we experimented with incorporating a large language model, FLAN-T5, as a post-processing 
step to improve the quality of the generated alt text. The idea was to take the raw output from the BLIP model 
and pass it through FLAN-T5 to refine the structure, improve clarity, and enhance readability. This approach 
leveraged the strong text understanding and generation capabilities of FLAN-T5 to correct grammatical issues, 
remove redundancies, and produce more polished descriptions.

Unlike the baseline, this introduced an additional layer of processing where the model operates purely on text 
embeddings rather than visual features. FLAN-T5 was prompted to rephrase or clean the generated captions, 
effectively acting as a text refinement module rather than a primary generator. This can be seen as a form 
of feature enhancement at the language level, where the semantic quality of the output is improved without modifying 
the underlying visual features.

\subsubsection{Alternative Model Testing}
To identify the best-performing model configuration, approximately 20 model variants were 
trained and evaluated across different combinations of architecture, size, and training data. 
The models tested included BLIP-small and BLIP-large, as well as Pix2Struct, each trained 
on one of three dataset configurations: our custom annotated physics dataset, the SciCap 
dataset (an open-source dataset of scientific figure captions), or a combination of both. 
This systematic comparison allowed us to evaluate the trade-off between model capacity, 
training data domain, and output quality. Key features such as image patch embeddings, 
vision encoder representations, and language decoder outputs were compared across runs 
using our internal evaluation rubric.


\section{Implementation}

This section describes the model architecture and implementation details across each revision, including training configurations, deployment decisions, and the key limitations observed at each stage.

\subsection{Revision 0 Pipeline}
In Revision 0, the pipeline consisted of using a pretrained BLIP model to generate alt text directly from input 
physics diagrams. Images from the dataset were passed into the model, which processed them through a vision encoder 
to extract visual features. These features were then fed into a language model to generate a corresponding textual description. 
The model used for this implementation was \texttt{Salesforce/blip-image-captioning-base}, which was selected for its 
strong baseline performance on image captioning tasks.

The model was fine-tuned for 3 epochs using the training configuration shown in Section 4. A relatively small batch size of 
1 was used due to hardware constraints, along with a learning rate of 0.00005. No advanced optimization techniques such as mixed precision or large-scale parallelization were applied in this revision. 
The training process focused on adapting the pretrained model to better handle physics-specific diagrams.

For deployment, the model was hosted using Hugging Face Inference, allowing for accessible and scalable testing. However, 
the inference time was approximately 2 minutes per image, which is relatively slow and not ideal for real-time applications.

Several limitations were observed in this baseline implementation. The generated alt text often lacked detail and structure, 
especially for complex STEM diagrams. In some cases, the model produced hallucinations, where it introduced incorrect or irrelevant 
information not present in the image. For example, diagrams involving forces or motion were sometimes described with incorrect labels or relationships, 
and technical concepts such as angular momentum were misrepresented or overly simplified. Additionally, the model struggled to capture deeper 
meaning or accurately describe relationships between components, resulting in descriptions that were either vague or partially incorrect.

Overall, while Revision 0 established a functional end-to-end pipeline for generating alt text, it highlighted key issues in accuracy, consistency, 
and inference efficiency that motivated further improvements in later revisions.

\subsection{Revision 1 Experiments}

\subsubsection{BLIP Re-training for Expanded Dataset} 
To expand the dataset, additional images were collected from an open-source physics textbook and preprocessed to 
remove irrelevant or low-quality images. 
Following this, the BLIP model was retrained multiple times on larger versions of the dataset.
We experimented with training configurations using 3 and 5 epochs, and evaluated performance across different categories of test images. 
We found that 3 epochs resulted in more accurate and consistent results. Ultimately, 3 epochs was chosen for training on the expanded dataset of 200 additional annotated images. 

\subsubsection{Category Routing System}

The category routing system was implemented as a preprocessing step in the pipeline. For each input image, the CLIP model first predicted the diagram category. This predicted category was then used to construct a category-specific prompt that guided the BLIP captioning model, enabling more structured and domain-specific alt text generation.

The addition of a subject field further enhanced this process by providing additional semantic context for prompt conditioning. This helped differentiate between diagrams with similar visual layouts but different conceptual meanings. By combining CLIP-based classification with prompt-guided BLIP caption generation, the system produced more consistent and relevant descriptions.

Overall, placing the CLIP classifier before the BLIP model introduced a lightweight routing mechanism that improved both prompt quality and caption accuracy while maintaining an efficient pipeline.

\subsubsection{LLM Polisher Post-Processing}
In Revision 1, we explored enhancing the pipeline by introducing a post-processing step using a large language model, specifically FLAN-T5, 
to improve the quality of the generated alt text. In this approach, the output from the BLIP model was passed into FLAN-T5, which was 
prompted to refine the text by improving sentence structure, clarity, and overall readability. This created a two-stage pipeline: 
image input (1) BLIP caption generation (2) FLAN-T5 text refinement.

Unlike the base pipeline, this additional step operated purely on textual embeddings, focusing on improving linguistic quality rather 
than extracting new visual features. The intention was to address issues observed in Revision 0, such as inconsistent formatting, 
repetition, and lack of fluency in the generated descriptions.

However, several limitations were identified during experimentation. FLAN-T5 often over-simplified or shortened the generated descriptions, 
removing important technical details that are critical for accurately describing STEM diagrams. In some cases, it cut sentences or omitted 
key relationships between components, reducing the usefulness of the alt text. Additionally, a broader limitation of using large language 
models is that many are either not fully open-source or require paid access for scalable usage, making them less practical for integration into a capstone-level system.

It is important to note that after experimenting with this, this approach was not included in the final implementation. Instead, a rule-based post-processing function was developed 
in the backend to manually clean and standardize the structure of the generated alt text. This method provided more control over formatting and 
ensured that important technical details were preserved, while still improving readability and consistency.

\subsubsection{Alternative Model Testing}
Each of the approximately 20 model variants was trained under consistent conditions to 
ensure fair comparison. Models were trained for 3 epochs, which was identified in earlier 
experiments as the optimal setting, using a batch size of 1 due to hardware constraints 
and a learning rate of 0.00005. For each configuration, outputs on a fixed set of test 
images were generated and scored using the internal evaluation rubric described in 
Section 6.1. Across all tested combinations, BLIP-small fine-tuned on our custom 
annotated physics dataset consistently produced the highest quality alt text, achieving 
the best balance of descriptive accuracy, structural consistency, and domain relevance. 
BLIP-large showed marginal improvement in some cases but was significantly slower to 
train and run inference. Pix2Struct variants underperformed across all dataset configurations, 
producing shorter, less accurate descriptions with a higher rate of hallucinated content. 
As a result, BLIP-small trained on the custom dataset was selected as the final model.

\section{Results and Evaluation}

This section presents the evaluation framework used to assess the quality of the generated alt text, along with results from both Revision 0 and Revision 1. Model outputs were evaluated by five internal evaluators using a pre-defined rubric that measures descriptive quality, accessibility, and educational value. Results are compared across revisions to quantify the impact of the improvements made.

\subsection{Evaluation Metrics} 
To evaluate the quality and effectiveness of the generated alt text, a combination of categorical, numerical, and qualitative metrics was used. 
These metrics were designed to assess not only the correctness of the descriptions but also their usability and impact for accessibility purposes.

The first metric, Sufficiency of Description, is a categorical measure rated on a scale from 1 to 3. It evaluates whether the generated alt text 
provides enough relevant information to convey the intended meaning of the diagram. A score of 3 is considered acceptable, indicating that 
the description is sufficiently detailed for user understanding.

The second metric, Length Appropriateness, is also a categorical measure on a 1 to 3 scale. This metric assesses whether the alt text is concise yet 
complete, avoiding descriptions that are either too brief or unnecessarily verbose. A score of 3 indicates that the description achieves a proper 
balance between completeness and conciseness.

The Accessibility and Usability metric is evaluated on a numerical scale from 0 to 3. This measures how well the generated alt text aligns with 
accessibility standards, particularly in terms of clarity and compatibility with assistive technologies such as screen readers. A score of 2 or higher 
is considered acceptable, reflecting that the output meets basic usability requirements aligned with WCAG 2.1 Level AA guidelines.

The fourth metric, Learning Impact, is also rated on a numerical scale from 0 to 3. This evaluates whether the generated alt text supports or enhances 
user understanding, particularly in educational contexts. A score of 2 or higher indicates that the description provides meaningful educational value.

Finally, Qualitative Feedback Notes are collected as free-form textual responses. This allows evaluators to provide detailed comments on aspects such 
as clarity, tone, structure, and any suggested improvements. These qualitative insights complement the quantitative metrics by capturing nuances that 
numerical scores may not fully reflect.

These metrics provide a comprehensive evaluation framework that balances technical accuracy, readability, accessibility, and educational 
usefulness, ensuring that the generated alt text meets the goals of the Reading4All project.

\subsection{Revision 0 Results}

Table~\ref{tab:rev0} presents the internal evaluation scores collected from five evaluators for the Revision 0 model outputs.

\begin{table*}[t]
\centering
\caption{Revision 0 Internal Evaluation Scores per Evaluator}
\label{tab:rev0}
\begin{tabular}{lcccc}
\toprule
\textbf{Evaluator} & \textbf{Sufficiency (1--3)} & \textbf{Length (1--3)} & \textbf{Accessibility (0--3)} & \textbf{Learning (0--3)} \\
\midrule
Person 1 & 2 & 3 & 3 & 1 \\
Person 2 & 2 & 3 & 3 & 2 \\
Person 3 & 2 & 3 & 3 & 1 \\
Person 4 & 2 & 2 & 3 & 1 \\
Person 5 & 2 & 2 & 3 & 1 \\
\midrule
\textbf{Mean}       & \textbf{2.0} & \textbf{2.6} & \textbf{3.0} & \textbf{1.2} \\
\textbf{Acceptable} & $\geq 3$     & $\geq 3$     & $\geq 2$     & $\geq 2$     \\
\bottomrule
\end{tabular}
\end{table*}

The results indicate that Accessibility/Usability consistently met the acceptable threshold 
with a perfect mean score of 3.0 across all evaluators. Length Appropriateness averaged 
2.6, falling slightly below the acceptable threshold of 3. Sufficiency of Description 
scored uniformly at 2.0, below the acceptable threshold of 3, indicating that the 
generated alt text did not yet convey sufficient information to fully capture diagram 
content. Learning Impact averaged 1.2, well below the acceptable threshold of 2, 
suggesting that the Revision 0 outputs provided limited educational value. These 
results highlighted the need for improvements in descriptive depth and domain-specific 
accuracy, motivating the experiments carried out in Revision 1.

\subsection{Revision 1 Comparisons}

Table~\ref{tab:rev1} presents the internal evaluation scores for the Revision 1 model outputs, 
collected using the same rubric and evaluator pool as Revision 0.

\begin{table*}[t]
\centering
\caption{Revision 1 Internal Evaluation Scores per Evaluator}
\label{tab:rev1}
\begin{tabular}{lcccc}
\toprule
\textbf{Evaluator} & \textbf{Sufficiency (1--3)} & \textbf{Length (1--3)} & \textbf{Accessibility (0--3)} & \textbf{Learning (0--3)} \\
\midrule
Person 1 & 2 & 3 & 3 & 2 \\
Person 2 & 3 & 3 & 3 & 2 \\
Person 3 & 2 & 3 & 3 & 2 \\
Person 4 & 2 & 3 & 3 & 1 \\
Person 5 & 3 & 3 & 3 & 2 \\
\midrule
\textbf{Mean}       & \textbf{2.4} & \textbf{3.0} & \textbf{3.0} & \textbf{1.8} \\
\textbf{Acceptable} & $\geq 3$     & $\geq 3$     & $\geq 2$     & $\geq 2$     \\
\bottomrule
\end{tabular}
\end{table*}

Compared to Revision 0, Revision 1 shows measurable improvement across all metrics. 
Length Appropriateness reached the acceptable threshold with a mean of 3.0, up from 
2.6, reflecting the impact of post-processing improvements on output structure and 
conciseness. Sufficiency of Description improved from 2.0 to 2.4, indicating that 
the expanded dataset and category routing system produced more informative descriptions, 
though it still falls short of the acceptable threshold of 3. Accessibility/Usability 
remained at 3.0, maintaining strong performance. Learning Impact showed the most 
notable relative gain, rising from 1.2 to 1.8, suggesting that the Revision 1 
improvements meaningfully increased the educational value of the generated alt text, 
though further improvement is still needed to reach the acceptable threshold of 2. 
Overall, Revision 1 demonstrates clear progress, with the most significant gains 
in output structure, length consistency, and learning support.

\section{Feedback and Plans}

This section reflects on the feedback received from the instructor, teaching assistant, and supervisor following Revision 0, and describes the plans and improvements carried out in response during Revision 1.

\subsection{Revision 0 Feedback} 
Feedback from the instructor, teaching assistant, and supervisor highlighted several key areas for improvement in Revision 0. 
Overall, the scope of focusing on physics diagrams was seen as appropriate; however, further narrowing of scope was suggested 
to improve model performance. A major point of feedback was that the quality of the generated alt text needed improvement, 
particularly in terms of clarity, flow, and consistency. Issues such as poor capitalization, awkward phrasing, and lack of structure were identified.

The teaching assistant specifically noted that some generated descriptions were unclear or did not fully make sense, indicating 
that the model struggled with accurately capturing relationships and meaning within complex diagrams. Additionally, it 
was suggested that the system should make it clearer that instructors can edit or refine the generated alt text before use, 
improving usability in real-world settings. From a system perspective, feedback also included improving clarity in features such as 
file handling and making system behaviors (e.g., history sorting) more transparent.

\subsection{Revision 1 Plan} 
Based on this feedback, our Revision 1 plan focused on conducting a series of targeted experiments to 
improve the quality, structure, and accuracy of the generated alt text. The work was divided into multiple 
subteams, each addressing a different aspect of the system.

The first area of focus was dataset expansion and restructuring, where the dataset was organized into 
physics-specific subsets (e.g., graphs, electric fields, block diagrams) and expanded to improve model generalization. 
A new subject field was added to support more structured training and enable future improvements.

The second component involved developing a routing and model selection system, where images would first 
be classified into a physics subset and then routed to a more specialized captioning approach. This aimed to move 
beyond a single generalized pipeline and instead tailor outputs based on diagram type.

Another key area was output improvement using LLMs, where we explored using models such as FLAN-T5 to refine 
the generated captions in terms of clarity, readability, and formatting. This was intended to directly address 
feedback regarding poor structure and flow in the generated alt text.

Finally, we conducted model experimentation and optimization, testing alternative models and configurations 
(e.g., different training epochs and architectures) to evaluate whether improvements over BLIP could be achieved. All experiments were documented and compared to guide final model selection.

\subsection{Revision 1 Improvement}
In Revision 1, several improvements were made to address the issues identified in earlier feedback. One of the most impactful changes was the 
introduction of a post-processing layer that cleans and standardizes the structure of the generated alt text. After BLIP generates the initial 
caption, a backend function is applied to correct formatting issues such as inconsistent capitalization, extra whitespace, repeated words, 
and unnecessary punctuation.

These improvements significantly enhanced the readability and professionalism of the generated descriptions without altering their original meaning. 
This was especially important for accessibility, as clearer and more structured text improves the experience for screen reader users.

In addition to structural improvements, minor enhancements were made to the quality of the alt text itself, resulting in slightly more detailed 
and coherent descriptions compared to Revision 0. While the core model remained similar, these refinements demonstrate how post-processing and data 
improvements can meaningfully improve output quality.

Revision 1 shows clear progress in addressing feedback by improving both the presentation and clarity of the generated alt text, while also laying 
the groundwork for further model and pipeline enhancements.


\section*{Team Contributions}

Each team member contributed to different aspects of the project across both revisions. The following summarizes the primary responsibilities of each contributor.

\begin{itemize}
    \item \textbf{Casey Francine Bulaclac:} BLIP training in Revision 0, image annotation in Revision 0, and LLM Polisher post-processing experimentation using FLAN-T5 in Revision 1.
    \item \textbf{Fiza Sehar:} BLIP research, training Revision 0 and Revision 1 models, dataset preprocessing, alt text annotation and refinement, CLIP-based category routing, prompt conditioning, and user testing analysis.
    \item \textbf{Dhruv Sardana:} Grammar and spelling review of the generated alt text outputs and project documentation to ensure clarity and correctness.
    \item \textbf{Nawaal Fatima:} Completed all outstanding report sections, including writing the Pix2Struct-related content, alternative model testing documentation, results comparisons, and section introductions.
    \item \textbf{Moly Mikhail:} Retraining the model on progressively larger datasets and comparing results across different epoch configurations. Also annotated 200 additional images for the expanded dataset.
\end{itemize}

% Custom bibliography entries only
\bibliographystyle{ieeetr}
\bibliography{custom}

\end{document}